\section{Zadanie 4}
Każdy z n graczy rzuca niezależnie od pozostałych kostką.
\begin{enumerate}[label=(\roman*)]
    \item\label{set04.task04.1} Za dowolną parę graczy którzy wyrzucili ten sam wynik grupa dostaje 1 punkt.
    Znaleźć wariancję liczby uzyskanych punktów.
    \item\label{set04.task04.2} Za dowolną parę graczy którzy wyrzucili ten sam wynik grupa dostaje tyle punktów
    ile wyrzucili. Znaleźć wariancję liczby uzyskanych punktów.
\end{enumerate}

\subsection*{Wstęp}
Niech $X_i$ będzie wynikiem rzutu $i$-tego gracza, dla $i = 1, \dots, n$.
Zmienne $X_i$ są niezależne i mają jednakowy rozkład: $P(X_i = k) = \frac{1}{6}$ dla $k \in \{1, 2, 3, 4, 5, 6\}$.

Będziemy korzystać ze wzoru na wariancję sumy zmiennych losowych:
$$\text{Var}\left(\sum_{k} A_k\right) = \sum_{k} \text{Var}(A_k) + \sum_{k \neq l} \text{Cov}(A_k, A_l)$$

\subsection*{Rozwiązanie podpunktu \ref{set04.task04.1}}

Definiujemy zmienną losową $Y$ jako całkowitą liczbę punktów.
Liczba punktów jest sumą punktów przyznanych każdej parze graczy. Jest $\binom{n}{2}$ takich par.

Dla każdej pary graczy $(i, j)$ z $1 \le i < j \le n$, zdefiniujmy indykator $Y_{ij}$:
$$Y_{ij} = \begin{cases} 1 & \text{gdy } X_i = X_j \\ 0 & \text{gdy } X_i \neq X_j \end{cases}$$

Całkowita liczba punktów $Y$ to oczywiście suma tych zmiennych:
$$Y = \sum_{1 \le i < j \le n} Y_{ij}$$

Aby znaleźć $\text{Var}(Y)$, musimy obliczyć $\text{Var}(Y_{ij})$ oraz $\text{Cov}(Y_{ij}, Y_{kl})$ dla $(i,j) \neq (k,l)$.
\subsubsection*{Wariancja}
Wartość oczekiwana: $E[Y_{ij}] = p = \frac{1}{6}$. \\
Wariancja: $\text{Var}(Y_{ij}) = p(1-p) = \frac{1}{6} \left(1 - \frac{1}{6}\right) = \frac{1}{6} \cdot \frac{5}{6} = \frac{5}{36}$.
\subsubsection*{Kowariancja}
Dla par rozłącznych Kowariancja jest równa 0 zastanówmy co dzieje się dla par mających jeden wspólny wskaźnik.
BSO weźmy pary $(i, j)$ i $(i, k)$, wtedy
$$P(Y_{ij} = 1\ \text{i}\ Y_{ik} = 1) = P(\text{wszyscy trzej gracze mają ten sam wynik}) = \frac{1}{36}$$
stąd
$$Cov(Y_{ij}, Y_{ik}) = \frac{1}{36} - (\frac{1}{6})^2 = 0 $$
zatem wszystkie kowariancje pomiedzy dwoma różnymi $Y_{ij}$ są równe $0$
\subsubsection{A zatem}
wariancja sumy jest po prostu sumą wariancji:
$$\text{Var}(Y) = \sum_{1 \le i < j \le n} \text{Var}(Y_{ij})$$
Liczba par $(i, j)$ wynosi $\binom{n}{2} = \frac{n(n-1)}{2}$.
$$\text{Var}(Y) = \binom{n}{2} \cdot \text{Var}(Y_{ij}) = \frac{n(n-1)}{2} \cdot \frac{5}{36} = \frac{5n(n-1)}{72}$$


\subsection*{Rozwiązanie podpunktu \ref{set04.task04.2}}
Definiujemy zmienną losową $Y$ jako całkowitą liczbę punktów.
Dla każdej pary graczy $(i, j)$ z $1 \le i < j \le n$, zdefiniujmy zmienną $Y_{ij}$:
$$Y_{ij} = \begin{cases} X_i & \text{gdy } X_i = X_j \\ 0 & \text{gdy } X_i \neq X_j \end{cases}$$

\subsubsection*{Moment pierwszego i drugiego rzędu}
\[
\E[Y_{ij}] = \sum_{k=1}^{6} k \cdot \frac{1}{36} = \frac{21}{36} = \frac{7}{12},
\]
\[
\E[Y_{ij}^2] = \sum_{k=1}^{6} k^2 \cdot \frac{1}{36} = \frac{91}{36},
\]
\[
Var(Y_{ij}) = \frac{91}{36} - \left(\frac{7}{12}\right)^2 = \frac{35}{16}.
\]

\subsubsection*{Kowariancja dwóch zmiennych}
Jeśli pary są rozłączne, kowariancja = 0.
Jeśli pary mają wspólnego gracza (\(Y_{ij}\) i \(Y_{ik}\)):
\[
\E[Y_{ij}Y_{ik}] = \sum_{k=1}^{6} k^2 \cdot \frac{1}{216} = \frac{91}{216},
\]
\[
Cov(Y_{ij}, Y_{ik}) = \frac{91}{216} - \left(\frac{7}{12}\right)^2 = \frac{35}{432}.
\]

\subsubsection*{Liczba par wskaźników ze wspólnym graczem}
Liczba uporządkowanych par o wspólnym wskaźniku (takich potrzebujemy do sumy) to: \\
\[
2 \cdot 3 \cdot \binom{n}{3} = n(n-1)(n-2).
\]

\subsubsection*{Wariancja całkowita}
Wiemy, że
\[
Y = \sum_{1 \le i < j \le n} Y_{ij}.
\]
Wtedy
\[
\begin{aligned}
\operatorname{Var}(Y)
&= \sum_{\{i,j\}} \operatorname{Var}(Y_{ij})
   + \sum_{\substack{(i < j) \neq (k < l)\\}}
     \operatorname{Cov}(Y_{ij}, Y_{ik})\\[6pt]
&= \binom{n}{2} \cdot \frac{35}{16} + n(n-1)(n-2) \cdot \frac{35}{432}\\[6pt]
&= \frac{35\,n(n-1)}{32} + \frac{35\,n(n-1)(n-2)}{432}\\[6pt]
&= 35 \cdot \frac{n(n-1)(2n+23)}{864}.
\end{aligned}
\]



\subsection*{Odpowiedź}
\[
\boxed{
    \text{(i)}\ \binom{n}{2} \cdot \frac{5}{36} = \frac{5n(n-1)}{72}
    \qquad
    \text{(ii)}\ \frac{35\,n(n-1)(2n+23)}{864}
    }
\]
